% ============================================================
%  Глава 9 — Сценарий «D метров вперёд»
% ============================================================
\chapter{Сценарий «D метров вперёд»}
\label{ch:scenario}

\section{Что мы строим}

Соберём всё: модель ($\mat{A}_d, \mat{B}_d$), регулятор ($\mat{K}_\text{first}$
из MPC), измерение состояния ($s, v, \theta, \omega, e_\text{int}$),
оркестратор (FSM). Получится \textbf{сценарий} --- от нажатия кнопки в UI
до полного отчёта.

\begin{ideabox}
Сценарий --- это «контракт» между UI и роботом. Он говорит:
\emph{при каких параметрах робот считается успешно справившимся,
а при каких --- нет}. В МПС-модуле сценарий один: «проехать $D$~м
вперёд» (с опциональным предварительным разворотом, о нём ---
глава~\ref{ch:turn_phase}). Каждый прогон --- финитный, не цикл.
\end{ideabox}

\section{Постановка задачи}

\textbf{Дано:}
\begin{itemize}
  \item Дистанция $D \in (0, D_\text{max}]$~м, по умолчанию $2.0$.
  \item Целевая скорость $v_\text{target} \in (0, v_\text{max}]$~м/с,
        по умолчанию $0.15$.
  \item Целевой курс $\varphi \in [-\pi, \pi]$ (см. главу~\ref{ch:picker}).
        По умолчанию $0$ --- едем прямо.
\end{itemize}

\textbf{Требуется:} от старта (положение $s_\text{start}$, курс
$\theta_\text{start}$) переместиться в точку с относительной
координатой $s = D$, курс удерживать у $\varphi$.

\textbf{Допустимые исходы:}
\begin{center}
\begin{tabular}{lll}
\toprule
\textbf{Событие} & \textbf{Условие} & \texttt{status}\\
\midrule
Достижение цели & $s \ge D - 0.05$~м & \texttt{reached}\\
Тайм-аут        & $t > 3\,D / v_\text{target}$ & \texttt{timeout}\\
Принудит. abort & MQTT \texttt{mps/scenario/abort} & \texttt{aborted}\\
Расходимость    & $\|\vect{x}\| > $ bound, NaN/Inf & \texttt{error}\\
\bottomrule
\end{tabular}
\end{center}

\section{Reference}

MPC требует $\vect{x}^\text{ref}_k$ --- куда хотим попасть на каждом
шаге. Самая простая стратегия: \textbf{линейный разгон до $D$ за
номинальное время}.

\begin{formulabox}[title={Reference сценария}]
\begin{equation}
  s_\text{ref}(k) = \min\big(D,\;k\,\Delta t \cdot v_\text{target}\big),
  \qquad
  \vect{x}^\text{ref}_k = \begin{pmatrix} s_\text{ref}(k)\\ v_\text{target}\\ \varphi\\ 0\\ 0 \end{pmatrix}.
  \label{eq:reference}
\end{equation}
\end{formulabox}

Геометрия: $s_\text{ref}$ растёт линейно с угловым коэффициентом
$v_\text{target}$, пока не достигнет $D$ --- дальше «плато». Курс
$\theta_\text{ref} = \varphi$ постоянен. Угловая скорость
$\omega_\text{ref} = 0$ --- робот не должен крутиться, удерживая курс.

\begin{notebox}
Почему $v_\text{ref} = v_\text{target}$ \emph{константа}, а не «по
правилу»? Потому что в постоянной части $s_\text{ref}$ (после
достижения плато) скорость и так логически 0, но \textbf{штраф $\mat{Q}$
по $v$ должен быть маленький} --- иначе MPC будет «тормозить»
заранее, не доезжая до $D$. В \filepath{config.yaml} мы видим
$Q_\text{diag}[v] = 1$ против $Q_\text{diag}[s] = 10$ --- скорость
важна сильно меньше позиции.
\end{notebox}

\section{Полный тик регулятора}

Один полный шаг \texttt{mps\_node} (50~Гц):

\begin{enumerate}
  \item \textbf{Прочитать измерение.} Берём последний \texttt{x\_meas} из
        одометрии. Конвертируем в относительные координаты
        \eqref{eq:relative_state}.
  \item \textbf{Построить reference.} По текущему $k$ собираем
        $\vect{x}^\text{ref}_k$ из \eqref{eq:reference}.
  \item \textbf{Шаг MPC.} \texttt{u = mpc.step(x\_meas, x\_ref)}.
        Под капотом --- $\vect{u}_0 = -\mat{K}_\text{first}(\vect{x} -
        \vect{x}^\text{ref})$ + клиппинг.
  \item \textbf{Sanity-check.} Если в $\vect{u}$ или $\vect{x}$ есть
        NaN/Inf --- немедленный \texttt{status='error'}.
  \item \textbf{Жёсткий cap угловой скорости в forward-режиме.}
        $|u_\omega| \le 0.5$~рад/с (а не $2.0$ как глобально) --- защита
        от случайного «винта» при удержании курса.
  \item \textbf{Опубликовать.} \texttt{cmd\_vel} в MQTT
        (\texttt{linear\_x}, \texttt{angular\_z}); телеметрию \texttt{x, u, y}
        --- в UI.
  \item \textbf{Проверить условия завершения.}
        Достигли? Тайм-аут? Watchdog (нет одометрии 3 тика подряд)?
\end{enumerate}

\begin{codebox}[title={\filepath{pi\_nodes/nodes/mps\_node.py:436}}]
\begin{lstlisting}[language=Python]
def _tick_drive(self, run, x):
    phi = run.target_heading
    # Reference: ramp s_ref to D, hold v_target, hold heading φ.
    s_ref = min(run.distance, run.drive_t * run.v_target)
    x_ref = np.array([s_ref, run.v_target, phi, 0.0, 0.0])

    try:
        u = self._mpc.step(x, x_ref=x_ref)
    except Exception as exc:
        self._finish_run('error', f'mpc.step: {exc}')
        return

    if not (np.all(np.isfinite(u)) and np.all(np.isfinite(x))):
        self._finish_run('error', 'NaN/Inf in u or x')
        return

    # Hard omega cap in forward scenario.
    u[1] = max(-self._omega_max_fwd,
               min(self._omega_max_fwd, u[1]))

    self._publish_cmd_and_telemetry(run, x, u)
    # ... check reached, timeout, watchdog
\end{lstlisting}
\end{codebox}

\section{Lifecycle прогона}

\begin{center}
\begin{tikzpicture}[
  node distance=12mm and 18mm,
  state/.style={
    draw, rounded corners=3pt, thick, fill=blue!8,
    align=center, font=\small, minimum width=3cm, minimum height=10mm,
  },
  flow/.style={-{Stealth[length=2.5mm]}, thick},
]
  \node[state] (idle) {\textbf{IDLE}};
  \node[state, right=of idle] (run) {\textbf{DRIVE\_FORWARD\_MPS}\\
                                     tick @ 50 Гц};
  \node[state, right=of run, fill=green!8]  (fin) {\textbf{FINISHED}\\
                                                   cmd\_vel = 0 $\times 3$};
  \node[below=14mm of run, font=\footnotesize, align=center]
       (events) {$s \ge D - 0.05$ (\texttt{reached})\\
                 $t > 3D/v$ (\texttt{timeout})\\
                 \texttt{abort}\\
                 NaN, $\|x\| > $bound (\texttt{error})};
  \draw[flow] (idle) -- node[above, font=\scriptsize]{mps/scenario/run} (run);
  \draw[flow, dashed] (run) -- (events);
  \draw[flow] (run) -- node[above, font=\scriptsize]{одно из событий} (fin);
  \draw[flow] (fin.south) .. controls +(0,-2.5) and +(0,-2.5) .. (idle.south);
\end{tikzpicture}
\end{center}

Failsafe в FINISHED: моторам отправляется команда $[0, 0]$
\textbf{три раза} подряд. Это страховка на случай, если первый пакет
потеряется в MQTT --- мотор-нода имеет встроенный watchdog
\texttt{cmd\_vel\_timeout}, который сам стопнет двигатели, но дублирование
не помешает.

\section{Метрики качества}

После окончания прогона на основе массива телеметрии считаются
\textbf{метрики}. Они отвечают на вопрос «как хорошо отработал
регулятор». Без них «работает» не значит «хорошо».

\begin{center}
\begin{tabular}{lll}
\toprule
\textbf{Метрика} & \textbf{Формула} & \textbf{Что значит}\\
\midrule
\texttt{overshoot} & $\max_t s(t) - D$, обрезано $\ge 0$ & «Пролетели мимо»\\
\texttt{settling\_time} & первое $t$, после которого $|s-D| < 0.02$ м & «Сколько успокаивались»\\
\texttt{control\_energy} & $\sum_k \vect{u}_k\T \mat{R}\vect{u}_k \Delta t$ & «Расход управления»\\
\texttt{ss\_error} & $|s(t_\text{end}) - D|$ & «Установившаяся ошибка»\\
\texttt{peak\_v} & $\max_t |v(t)|$ & «Пиковая скорость»\\
\texttt{peak\_omega} & $\max_t |\omega(t)|$ & «Пиковая угловая скорость»\\
\bottomrule
\end{tabular}
\end{center}

\begin{ideabox}
\textbf{Учебный смысл метрик.} Можно специально подобрать «плохой»
$\mat{Q}$ или $\mat{R}$ так, чтобы:
\begin{itemize}
  \item получить большой \texttt{overshoot} (например, $r_v$ малая ---
        регулятор «топит педаль»);
  \item получить большой \texttt{settling\_time} ($r_v$ большая ---
        регулятор «медленный»);
  \item получить большую \texttt{ss\_error} ($q_s$ малая --- регулятору
        не критично, доехал ли).
\end{itemize}
Это и есть «учебный пульт» из спеки: пробуем разные веса, смотрим, что
получилось. Курсовая.
\end{ideabox}

\section{Safety}

Чтобы случайный «учебный пульт» не сломал железо, есть многоуровневая
защита:

\begin{enumerate}
  \item \textbf{Pre-validate}: запрос с $D > D_\text{max}$ или
        $v_\text{target} > v_\text{max}$ отбрасывается ещё на REST
        (HTTP 400) и ещё раз на Pi (\texttt{mps/error},
        \texttt{error\_type="precondition"}).
  \item \textbf{Watchdog по одометрии}: если 3 тика подряд нет
        \texttt{odom} --- прогон завершается ошибкой
        \texttt{watchdog}.
  \item \textbf{FSM-lock}: в \texttt{DRIVE\_FORWARD\_MPS} игнорируются
        voice/joystick/детекции мяча. Только \texttt{abort} переключает
        состояние.
  \item \textbf{Жёсткий cap $u_\omega$}: в forward-режиме $|u_\omega|
        \le 0.5$~рад/с (это меньше, чем общий $u_\text{max}[1] = 2$).
  \item \textbf{Failsafe}: при любом FINISHED тройной \texttt{cmd\_vel = 0}.
\end{enumerate}

\section{Симулятор}

Параллельно с роботом существует \emph{идеальный} симулятор ---
\filepath{compute\_node/mps\_runner.py}. Он гоняет ровно ту же
последовательность \texttt{mpc.step(x, x\_ref)}, но вместо одометрии
обновляет $\vect{x}$ по дискретной модели:

\[
  \vect{x}_{k+1} = \mat{A}_d\,\vect{x}_k + \mat{B}_d\,\vect{u}_k.
\]

Без шумов, без motor lag, без проскальзывания. Если запустить «Run on Sim»
из UI, прогон выполнится за $\sim 100$~мс прямо на бэкенде --- никакого
железа. Это нужно:

\begin{itemize}
  \item для UI --- быстрый итеративный подбор $\mat{Q}, \mat{R}$ без
        возни с роботом;
  \item для тестов --- юнит-тесты не требуют железа;
  \item для сравнения --- разница между sim и реальным прогоном
        показывает, насколько модель отстаёт от реальности.
\end{itemize}

\section{Что мы получили}

\begin{itemize}
  \item Сценарий определяет, \emph{что значит} «успешный прогон»:
        \texttt{reached}, \texttt{timeout}, \texttt{aborted}, \texttt{error}.
  \item Reference --- линейный разгон $s_\text{ref}$ до $D$ за
        номинальное время.
  \item Один тик --- 7 шагов от чтения одометрии до публикации
        \texttt{cmd\_vel}.
  \item Метрики --- объективная мера качества, считаются по концу
        прогона.
  \item Безопасность --- 5 уровней защиты от ломки железа учебным
        пультом.
  \item Симулятор --- идеальная модель, быстрая проверка без робота.
\end{itemize}

Это была \emph{одно}фазная картина. Но в текущей ветке мы добавили
\textbf{предварительный разворот} к выбранному курсу. Этим и
займёмся.
